\chapter{Introducción}
% cambiar referencias con \cite{xd} 
La visión desempeña un papel crucial en la vida diaria, al ser el principal medio a través del cual percibimos e interactuamos con nuestro entorno. La visión, el más dominante de nuestros sentidos, desempeña un papel fundamental en cada faceta y etapa de nuestras vidas \cite{OMS_ceguera}. Según la OMS, al menos 2.2 mil millones de personas en el mundo tienen alguna deficiencia visual, y de estas, al menos mil millones podrían haberse evitado o aún no se han tratado \cite{OMS_ceguera}.

Las personas con discapacidad visual enfrentan desafíos diarios que limitan su independencia y acceso a diversas actividades cotidianas como la movilidad y la navegación, incluso ponen en riesgo su seguridad en entornos desconocidos. Esta discapacidad, puede variar en grado, desde dificultades para ver objetos cercanos, objetos lejanos, hasta la ceguera total. Las barreras físicas y sociales que encuentran a menudo les impiden desenvolverse con facilidad en entornos urbanos, lo que puede afectar su bienestar emocional y psicológico.

La creación de entornos accesibles es fundamental para mejorar la calidad de vida de las personas con discapacidad visual.  Las calles y las aceras pueden ser difíciles de transitar sin la ayuda de un bastón o un perro guía. Además, las personas con discapacidad visual pueden tener dificultades para subir escaleras o usar medios de transporte público que no están adaptados para ellos \cite{Covarrubias2025_desafios}.

Se ha señalado que la movilidad y la orientación son preocupaciones constantes para las personas ciegas \cite{PoonaSchool_mob}. Desplazarse por entornos desconocidos o altamente transitados, identificar obstáculos o acceder al transporte público representa un gran desafío. Las calles concurridas, los edificios sin señalización adecuada o las superficies irregulares suelen convertirse en barreras que limitan su independencia \cite{PoonaSchool_mob}.

En los últimos años, se ha observado un interés creciente por el desarrollo de tecnologías inclusivas que favorezcan la autonomía de personas con discapacidad visual. Este proyecto surge a partir de esa necesidad, con el objetivo de explorar nuevas formas de asistencia en tiempo real que permitan mejorar la orientación espacial y asegurar la seguridad de los usuarios en caso de estar expuestos a siniestros como lo son los incendios.

Se han propuesto diversas soluciones para asistir la movilidad de personas con discapacidad visual. A nivel nacional, el IPN reportó un prototipo portátil basado en sensores ultrasónicos que coloca dispositivos en el cuello, rodillas y abdomen para detectar obstáculos y avisar al usuario mediante vibraciones y alertas de voz \cite{IPN_invidentes}. Asimismo, estudiantes del IPN desarrollaron la “Primera Guía Robot” (un pequeño robot asistente) que integra cámaras 3D, algoritmos de IA y retroalimentación auditiva/vibratoria para guiar al usuario en tiempo real evitando obstáculos \cite{AhorraSeguros2025_guia}. 

\begin{itemize}
    \item Lentes de Meta Ray-Ban: Las gafas inteligentes Ray-Ban Meta. Aunque inicialmente no fueron diseñadas específicamente para personas con discapacidad visual, sus capacidades impulsadas por inteligencia artificial ofrecen funcionalidades de asistencia. La función de Meta AI con cámara permite que las gafas describan el entorno en tiempo real, identifiquen objetos, traduzcan texto y respondan a preguntas basadas en lo que ven. \cite{meta_rayban_2024}

    \item Lentes inteligentes con IA: En Cinvestav Guadalajara crearon el prototipo SmartOjo (2019), unas gafas equipadas con doble cámara, parlantes y una microcomputadora que procesan el entorno en tiempo real. SmartOjo permite navegar, evadir obstáculos móviles, reconocer rostros, leer texto y billetes, entregando instrucciones auditivas al usuario \cite{cinestavlentes}. Similarmente, la UPC coordinó el proyecto BCNMob (2024) usando las gafas Biel Glasses (tecnología 5G + realidad extendida): estas gafas 3D detectan y evitan obstáculos urbanos y ofrecen al usuario indicaciones visuales y sonoras personalizadas mediante IA y posicionamiento Galileo \cite{CIT2024_gafas}. Recientemente, el reto MIT Solve 2023 presentó el proyecto Eye in AI, un par de gafas inteligentes que mediante redes neuronales describen el entorno en voz alta (reconocimiento de objetos, rostros, texto a voz) \cite{eyeAI2023}.

    \item Aplicaciones móviles: Existen apps de navegación accesible por voz. Como, Lazarillo App, en México, es un GPS accesible que da instrucciones de audio en tiempo real al desplazarse por la ciudad o en interiores, permitiendo explorar entornos de forma más segura \cite{Lazarillo2023_app}. También Google Lookout y Seeing AI de Microsoft ofrecen reconocimiento de objetos y lectura de texto mediante IA, aunque su enfoque principal es asistencia visual, no rutas específicas.

    \item Tecnologías emergentes: Se han investigado sistemas con realidad virtual/aumentada y feedback háptico. El estudio “Nature Machine Intelligence 2025” describió un sistema chino compuesto por gafas con cámara 3D, un ‘A-skin’ vibratorio en la muñeca y audio óseo. Este dispositivo emite señales hápticas direccionales y sonidos envolventes discretos para guiar al usuario. \cite{Bordonaba2025_navegador}

    \item Dispositivos comerciales: En el mercado existen ayudas consolidadas, aunque no todas enfocadas en navegación. Por ejemplo, OrCam MyEye y eSight son gafas inteligentes que identifican texto y objetos para ciegos; el bastón “smart” WeWALK incorpora sensores ultrasónicos, GPS y conectividad móvil para dar alertas de voz. Sin embargo, muchos de estos productos no proveen navegación guiada activa.
\end{itemize}

La mayoría de los trabajos académicos se centran en prototipos experimentales con visión computacional y retroalimentación sonora para evitar obstáculos y mejorar la orientación. En este trabajo terminal se propone integrar lentes adaptado con IA y un sistema de rutas en tiempo real. 

Para desarrollar el proyecto, resulta esencial examinar la problemática que da origen a esta investigación. Comprender sus causas y consecuencias permite justificar la pertinencia del proyecto y definir el enfoque adoptado. En la siguiente sección se aborda el planteamiento del problema.